\chapter{Of Supply, Hydroponics, and the Ration}
\label{art:supply}

\articlenote{In which is set down the charge of Supply: the growing of food, the keeping of its stores and the silo's stores of water and air, and the setting of the ration by which every citizen's share is known and kept fair.}

\section{The Charge of Supply}

Supply is charged with the growing of food, the keeping of the silo's stores of food, water, and air, and the setting and distribution of the ration by which those stores are shared among the citizens.

\begin{clauses}
\item The Department is led by a Head of Supply, raised up from among its own trades by the shadowing entered into under Article 13, and confirmed in office by the Mayor.
\item Every citizen assigned to a trade of Supply serves under the Head of Supply in the ordinary conduct of that trade, and remains bound in every other respect to this Pact, to the Sheriff, and to Judicial, exactly as any other citizen is bound.
\end{clauses}

\section{Hydroponics and the Keeping of Livestock}

The growing of food is carried out in the hydroponic farms of the Mids, as Article 1 sets down, and in such further plots and pens as Supply is given charge of.

\begin{clauses}
\item Supply may keep livestock for the ration's sake, in such number and kind as the silo's stores of feed and space allow, and no citizen not assigned to that keeping may tend, slaughter, or dispose of an animal kept under this clause.
\item A citizen who tends the farms or the livestock of Supply is answerable for the yield of that tending, and a citizen found to have wasted, spoiled, or diverted food or feed from its proper use is answerable under Article 17.
\end{clauses}

\section{The Ration}

The ration is set by the Head of Supply from the full count of the rolls under Article 2, and no other number, that every citizen's share may be known and kept fair.

\begin{clauses}
\item A citizen's ration is theirs by right from naming to death, save where this Pact provides otherwise, and is not increased for labor, office, or seniority, though it may be adjusted for age, health, or condition where Article 13 or Article 19 expressly provides.
\item Where the silo's stores fall short of the ration this Pact provides, the Head of Supply shall report the shortfall to the Mayor and to Judicial without delay, and the ration shall be lessened by the smallest measure the shortfall requires, and lessened alike for every citizen, without exception of floor, trade, or office.
\item Goods too many or too heavy for a citizen's own carrying shall pass between the Down Deep and the floors of Supply by the freight lift under Article 9, and from Supply to every floor by such further means as this Pact elsewhere provides.
\end{clauses}

\section{Stores of Water and Air}

Supply shall keep such stores of water and air as Mechanical delivers under Article 9, and shall distribute them by the same ration this Article sets down.

\begin{clauses}
\item Where the works of water or air under Article 9 fail or fall short, Supply shall draw upon its stores to hold the ration steady for so long as those stores allow, and shall report to the Head of Mechanical and to the Mayor the day those stores will not.
\end{clauses}

\section{Sanctioned Animals Kept by Citizens}

A citizen who wishes to keep an animal for companionship, and not for the ration, may petition Supply to sanction it; Supply shall weigh the silo's stores of feed against the asking, and may grant, limit, or refuse the sanction as those stores require.

\begin{clauses}
\item An animal kept under this section draws upon the household's own ration for its feed, and is counted, for every purpose of this Pact, as the sanctioning citizen's to answer for.
\item An animal kept without sanction under this section is treated as livestock diverted from Supply under Section 2, and the citizen keeping it is answerable under Article 17.
\end{clauses}
